\documentclass[12pt,a4paper]{article}
% =============================================================================
% housestyle.tex — CIE A Level Physics Revision Booklets
% House style preamble: input this file at the top of every topic booklet
% =============================================================================

% ── Core packages ─────────────────────────────────────────────────────────────
\usepackage[a4paper, top=2.2cm, bottom=2.2cm, left=2.2cm, right=2.2cm]{geometry}
\usepackage{amsmath, amssymb}
\usepackage{booktabs}
\usepackage{array}
\usepackage{xcolor}
\usepackage[most]{tcolorbox}
\usepackage{enumitem}
\usepackage{hyperref}
\usepackage{microtype}
\usepackage{fancyhdr}
\usepackage{titlesec}
\usepackage{parskip}
\usepackage{multicol}
\usepackage{tocloft}

% ── Colour palette ────────────────────────────────────────────────────────────
\definecolor{cieprimary}{HTML}{003366}      % dark blue — headings, banner
\definecolor{ciepink}{HTML}{F77E9D}         % mid pink — formula frame, accent rule
\definecolor{ciepinkbg}{HTML}{FDE5EB}       % pale pink — formula box fill
\definecolor{ciedefnbg}{HTML}{DCE8FA}       % light blue — definition box fill
\definecolor{cietipbg}{HTML}{DCF5E6}        % pale green — exam tip fill
\definecolor{cietipframe}{HTML}{006040}     % dark green — exam tip frame
\definecolor{ciewarnbg}{HTML}{FFE6E6}       % pale red — common mistake fill
\definecolor{ciewarnframe}{HTML}{961E1E}    % dark red — common mistake frame
\definecolor{cieexbg}{HTML}{F5F5F5}         % light grey — worked example fill
\definecolor{cieexframe}{HTML}{AAAAAA}      % mid grey — worked example frame

% ── Page style ────────────────────────────────────────────────────────────────
\setlength{\headheight}{15pt}
\pagestyle{fancy}
\fancyhf{}
\fancyhead[L]{\small\textcolor{cieprimary}{\textbf{CIE A Level Physics}}}
\fancyhead[R]{\small\textcolor{cieprimary}{\bookletheader}}   % defined per topic
\fancyfoot[C]{\small\thepage}
\renewcommand{\headrulewidth}{0.5pt}
\renewcommand{\headrule}{\hbox to\headwidth{\color{cieprimary}\leaders\hrule height \headrulewidth\hfill}}

% ── Section / subsection formatting ──────────────────────────────────────────
% Section: dark blue, bold, with a pink rule beneath
\titleformat{\section}
  {\large\bfseries\color{cieprimary}}
  {\thesection}{0.6em}{}
  [\vspace{2pt}{\color{ciepink}\titlerule[1.5pt]}]

\titleformat{\subsection}
  {\normalsize\bfseries\color{cieprimary}}
  {\thesubsection}{0.5em}{}

\titleformat{\subsubsection}
  {\normalsize\bfseries}
  {}{0pt}{}

% ── TOC formatting (tocloft) ──────────────────────────────────────────────────
\renewcommand{\cftsecfont}{\bfseries\color{cieprimary}}
\renewcommand{\cftsecpagefont}{\bfseries\color{cieprimary}}
\renewcommand{\cftsubsecfont}{\color{cieprimary}}
\renewcommand{\cftsubsecpagefont}{\color{cieprimary}}
\setlength{\cftbeforesecskip}{4pt}

% ── tcolorbox environments ────────────────────────────────────────────────────
\tcbuselibrary{skins, breakable}

% Definition box — blue
\newtcolorbox{defnbox}[1]{
  enhanced, breakable,
  colback=ciedefnbg, colframe=cieprimary,
  fonttitle=\bfseries\color{white},
  title={Definition: #1},
  attach boxed title to top left={yshift=-2mm, xshift=4mm},
  boxed title style={colback=cieprimary, colframe=cieprimary, arc=2pt},
  arc=3pt, boxrule=1pt,
  top=6pt
}

% Formula box — pink
\newtcolorbox{formulabox}[1]{
  enhanced, breakable,
  colback=ciepinkbg, colframe=ciepink,
  fonttitle=\bfseries\color{cieprimary},
  title={#1},
  attach boxed title to top left={yshift=-2mm, xshift=4mm},
  boxed title style={colback=ciepink, colframe=ciepink, arc=2pt,
                     fontupper=\bfseries\color{white}},
  arc=3pt, boxrule=1pt,
  top=6pt
}

% Exam tip box — green
\newtcolorbox{tipbox}{
  enhanced, breakable,
  colback=cietipbg, colframe=cietipframe,
  fonttitle=\bfseries\color{white},
  title={$\checkmark$\; Exam Tip},
  attach boxed title to top left={yshift=-2mm, xshift=4mm},
  boxed title style={colback=cietipframe, colframe=cietipframe, arc=2pt},
  arc=3pt, boxrule=1pt,
  top=6pt
}

% Common mistake box — red
\newtcolorbox{warnbox}{
  enhanced, breakable,
  colback=ciewarnbg, colframe=ciewarnframe,
  fonttitle=\bfseries\color{white},
  title={!\; Common Mistake},
  attach boxed title to top left={yshift=-2mm, xshift=4mm},
  boxed title style={colback=ciewarnframe, colframe=ciewarnframe, arc=2pt},
  arc=3pt, boxrule=1pt,
  top=6pt
}

% Worked example box — grey
\newtcolorbox{examplebox}[1]{
  enhanced, breakable,
  colback=cieexbg, colframe=cieexframe,
  fonttitle=\bfseries\color{cieprimary},
  title={Example #1},
  attach boxed title to top left={yshift=-2mm, xshift=4mm},
  boxed title style={colback=cieexframe, colframe=cieexframe, arc=2pt,
                     fontupper=\bfseries\color{white}},
  arc=3pt, boxrule=1pt,
  top=6pt
}

% Formula summary box — pink frame, used around the summary table
\newtcolorbox{summarybox}{
  enhanced, breakable,
  colback=ciepinkbg, colframe=ciepink,
  fonttitle=\bfseries\color{white},
  title={Formula Summary Sheet},
  attach boxed title to top left={yshift=-2mm, xshift=4mm},
  boxed title style={colback=ciepink, colframe=ciepink, arc=2pt},
  arc=3pt, boxrule=1.5pt,
  top=8pt
}

% Mark scheme box — subtle blue tint
\newtcolorbox{markbox}{
  enhanced, breakable,
  colback=ciedefnbg!40, colframe=cieprimary!60,
  arc=3pt, boxrule=0.8pt,
  top=4pt, bottom=4pt
}

% Checklist box — pink accent
\newtcolorbox{checklistbox}{
  enhanced, breakable,
  colback=ciepinkbg!50, colframe=ciepink,
  fonttitle=\bfseries\color{white},
  title={Revision Checklist},
  attach boxed title to top left={yshift=-2mm, xshift=4mm},
  boxed title style={colback=ciepink, colframe=ciepink, arc=2pt},
  arc=3pt, boxrule=1.5pt,
  top=8pt
}

% ── Hyperlinks ────────────────────────────────────────────────────────────────
\hypersetup{
  colorlinks=true,
  linkcolor=cieprimary,
  urlcolor=cieprimary
}

% ── Utility: mark allocation right-aligned ────────────────────────────────────
% Usage: \markalloc{2}  produces a right-aligned [2 marks]
\newcommand{\markalloc}[1]{\hfill\textit{[#1 marks]}}

% ── Title page command ────────────────────────────────────────────────────────
% Usage: \maketitlepage{Topic 13}{Gravitational Fields}{sub-line}{bullet list body}
\newcommand{\maketitlepage}[4]{%
  \begin{titlepage}
  \thispagestyle{empty}
  % Banner
  \noindent\colorbox{cieprimary}{%
    \begin{minipage}{\textwidth}
      \vspace{10pt}
      \centering
      {\fontsize{16}{20}\selectfont\bfseries\color{white}
        CIE International A Level Physics (9702)}\\[6pt]
      {\fontsize{22}{28}\selectfont\bfseries\color{white} #1}\\[4pt]
      {\fontsize{26}{32}\selectfont\bfseries\color{white} #2}\\[6pt]
      {\fontsize{13}{18}\selectfont\color{white!80!cieprimary} #3}\\[10pt]
    \end{minipage}%
  }
  % Pink accent rule
  \noindent{\color{ciepink}\rule{\textwidth}{4pt}}
  \vspace{1.2cm}
  \begin{center}
    {\large\bfseries\color{cieprimary} Revision Booklet}\\[0.3cm]
    {\normalsize Syllabus 9702 \quad|\quad A Level (A2) \quad|\quad 2025--2027}
  \end{center}
  \vspace{0.8cm}
  \noindent\colorbox{ciepinkbg}{%
    \begin{minipage}{0.96\textwidth}
      \vspace{6pt}
      {\bfseries\color{cieprimary} This booklet covers:}
      \vspace{2pt}
      #4
      \vspace{4pt}
    \end{minipage}%
  }
  \end{titlepage}%
}
% =============================================================================


% ── Topic-specific header (used in page headers) ──────────────────────────────
\newcommand{\bookletheader}{Topics 14, 15 \& 16 — Thermal Physics}

% =============================================================================
\begin{document}
% =============================================================================

% ── Title page ────────────────────────────────────────────────────────────────
\maketitlepage
  {Topics 14, 15 \& 16}
  {Thermal Physics}
  {Temperature \(\cdot\) Ideal Gases \(\cdot\) Thermodynamics}
  {%
    \begin{itemize}[leftmargin=1cm, itemsep=1pt, label=\textbullet]
      \item Thermal Equilibrium and Temperature Scales
      \item Specific Heat Capacity and Specific Latent Heat
      \item The Mole and the Ideal Gas Equation
      \item Kinetic Theory of Gases
      \item Internal Energy
      \item The First Law of Thermodynamics
      \item Formula Summary, Worked Examples, Exam Practice \& Checklist
    \end{itemize}
  }

% ── Table of contents ─────────────────────────────────────────────────────────
\tableofcontents
\newpage

% =============================================================================
\section{Core Concepts and Definitions}
% =============================================================================

\subsection{What is Temperature?}

Temperature is a measure of the \textbf{average random kinetic energy} of the
molecules in a substance. It is \emph{not} the same as thermal (internal) energy: a
large cold object may contain more internal energy than a small hot object, yet have a
lower temperature.

\begin{defnbox}{Thermal Equilibrium}
Two objects are in \textbf{thermal equilibrium} when there is no net transfer of energy
between them — they are at the same temperature.
\end{defnbox}

This is the basis of the \textit{zeroth law of thermodynamics}: if objects A and B are
each in thermal equilibrium with object C, then A and B are also in thermal equilibrium
with each other.

% =============================================================================
\section{Topic 14 — Temperature}
% =============================================================================

\subsection{Temperature Scales}

\subsubsection{The Celsius Scale}
The Celsius scale is defined by two fixed points:
\begin{itemize}
  \item $0\,{}^{\circ}\text{C}$ — the ice point (pure water–ice equilibrium at $1\,\text{atm}$)
  \item $100\,{}^{\circ}\text{C}$ — the steam point (pure water–steam equilibrium at $1\,\text{atm}$)
\end{itemize}

\subsubsection{The Thermodynamic (Kelvin) Scale}
The thermodynamic scale does \textbf{not} depend on the properties of any particular
substance. Its zero point, \textbf{absolute zero}, is the temperature at which molecules
have the minimum possible kinetic energy. No substance can be cooled below this point.

\begin{formulabox}{Temperature Conversion}
\begin{equation}
  T\,/\,\mathrm{K} = \theta\,/\,{}^{\circ}\mathrm{C} + 273.15
\end{equation}
\begin{tabular}{ll}
$T$ & = thermodynamic temperature (K) \\
$\theta$ & = Celsius temperature (${}^{\circ}$C)
\end{tabular}
\end{formulabox}

\begin{tipbox}
Always convert temperatures to kelvin before substituting into gas-law equations.
A temperature \emph{difference} $\Delta\theta$ in ${}^{\circ}$C is numerically equal to
$\Delta T$ in K, so conversion is not needed for differences — but always required for
absolute values in $pV = nRT$.
\end{tipbox}

\subsection{Thermometric Properties}

A \textbf{thermometric property} is a physical quantity that varies measurably and
reproducibly with temperature, and can therefore be used to measure it.

\medskip
\begin{center}
\renewcommand{\arraystretch}{1.3}
\begin{tabular}{ll}
\toprule
\textbf{Thermometric property} & \textbf{Type of thermometer} \\
\midrule
Volume of a gas at constant pressure & Gas thermometer \\
Resistance of a metal (e.g.\ platinum) & Resistance thermometer (PRT) \\
E.m.f.\ of a thermocouple & Thermocouple thermometer \\
Density of a liquid & Liquid-in-glass thermometer \\
\bottomrule
\end{tabular}
\end{center}

\medskip
Different thermometers may give slightly different readings between the fixed points
because their thermometric properties do not all vary linearly with temperature in
identical ways. The thermodynamic scale avoids this problem entirely.

\subsection{Specific Heat Capacity}

\begin{defnbox}{Specific Heat Capacity $c$}
The \textbf{specific heat capacity} of a substance is the energy required to raise the
temperature of $1\,\mathrm{kg}$ of that substance by $1\,\mathrm{K}$
(or $1\,{}^{\circ}\mathrm{C}$).
\end{defnbox}

\begin{formulabox}{Energy and Temperature Change}
\begin{equation}
  Q = mc\Delta T
\end{equation}
\begin{tabular}{ll}
$Q$        & = energy supplied or removed (J) \\
$m$        & = mass (kg) \\
$c$        & = specific heat capacity (J\,kg$^{-1}$\,K$^{-1}$) \\
$\Delta T$ & = temperature change (K or ${}^{\circ}$C) \\
\end{tabular}
\end{formulabox}

\begin{tipbox}
In calorimetry experiments the energy supplied by an electric heater is
$Q = IVt$. Setting $IVt = mc\Delta T$ and solving for $c$ is a standard exam
calculation.
\end{tipbox}

\subsection{Specific Latent Heat}

During a \textbf{change of state} the temperature remains constant while energy is
supplied or removed. This energy changes the \emph{potential} energy of the molecules
(breaking or forming bonds) without changing their kinetic energy.

\begin{defnbox}{Specific Latent Heat $L$}
The \textbf{specific latent heat} of a substance is the energy required to change the
state of $1\,\mathrm{kg}$ of the substance at constant temperature.
\end{defnbox}

\begin{formulabox}{Latent Heat}
\begin{equation}
  Q = mL
\end{equation}
\begin{tabular}{ll}
$Q$ & = energy supplied or removed (J) \\
$m$ & = mass (kg) \\
$L$ & = specific latent heat (J\,kg$^{-1}$) \\
\end{tabular}
\end{formulabox}

Two types must be distinguished:
\begin{itemize}
  \item \textbf{Specific latent heat of fusion} $L_f$ — solid $\leftrightarrow$ liquid
  \item \textbf{Specific latent heat of vaporisation} $L_v$ — liquid $\leftrightarrow$ gas
\end{itemize}

In general $L_v \gg L_f$ because during vaporisation molecules must be completely
separated, requiring far more work against intermolecular forces than during melting.

\begin{warnbox}
Do not apply $Q = mc\Delta T$ when a substance is changing state. During a change of
state the temperature does \emph{not} change — use $Q = mL$ instead. A common error is
forgetting to split a cooling/heating calculation into separate stages: warming,
change of state, then further warming.
\end{warnbox}

% =============================================================================
\section{Topic 15 — Ideal Gases}
% =============================================================================

\subsection{The Mole and Avogadro's Constant}

\begin{defnbox}{The Mole}
One \textbf{mole} of a substance contains exactly $N_A = 6.02 \times 10^{23}$
particles, where $N_A$ is the \textbf{Avogadro constant}.
\end{defnbox}

Key relationships:
\begin{align}
  N &= nN_A \qquad \text{(number of molecules from moles)} \\
  n &= \frac{m_{\text{sample}}}{M_r} \qquad \text{(moles from mass and molar mass)}
\end{align}

\subsection{The Ideal Gas Equation}

The empirical gas laws (Boyle's, Charles's, Pressure Law) combine into a single
equation of state for an ideal gas:

\begin{formulabox}{Ideal Gas Equation}
\begin{align}
  pV &= nRT \\[4pt]
  pV &= NkT
\end{align}
\begin{tabular}{ll}
$p$ & = pressure (Pa) \\
$V$ & = volume (m$^3$) \\
$n$ & = amount of substance (mol) \\
$R$ & = molar gas constant $= 8.31\,\mathrm{J\,mol^{-1}\,K^{-1}}$ \\
$T$ & = thermodynamic temperature (K) \\
$N$ & = number of molecules \\
$k$ & = Boltzmann constant $= 1.38 \times 10^{-23}\,\mathrm{J\,K^{-1}}$ \\
\end{tabular}
\end{formulabox}

The two constants are related by:
\begin{equation}
  k = \frac{R}{N_A}
\end{equation}

A gas \textbf{obeys} the ideal gas equation if $pV \propto T$ — equivalently, if
$pV/T$ is constant for a fixed mass of gas.

\begin{tipbox}
For a fixed mass of gas changing state:
\[
  \frac{p_1 V_1}{T_1} = \frac{p_2 V_2}{T_2}
\]
Use $pV = nRT$ when moles are given; use $pV = NkT$ when number of molecules is given.
\end{tipbox}

\subsection{Kinetic Theory of Gases}

\subsubsection{Assumptions of the Kinetic Model}

\begin{itemize}
  \item The gas consists of a \textbf{large number} of identical molecules.
  \item Molecules are in \textbf{continuous, random motion}.
  \item The \textbf{volume of the molecules} is negligible compared to the volume of
        the container.
  \item \textbf{Intermolecular forces} are negligible except during collisions.
  \item All collisions (molecule–molecule and molecule–wall) are \textbf{perfectly
        elastic}.
  \item The \textbf{duration of a collision} is negligible compared to the time
        between collisions.
\end{itemize}

\subsubsection{Pressure from Kinetic Theory}

By considering molecules bouncing elastically off the walls of a container and
extending from one dimension to three, it can be shown that:

\begin{formulabox}{Kinetic Theory Equation}
\begin{equation}
  pV = \tfrac{1}{3}\,Nm\langle c^2 \rangle
\end{equation}
\begin{tabular}{ll}
$N$                     & = total number of molecules \\
$m$                     & = mass of one molecule (kg) \\
$\langle c^2 \rangle$   & = mean-square speed (m$^2$\,s$^{-2}$) \\
\end{tabular}
\end{formulabox}

The \textbf{root-mean-square speed} is defined as:
\begin{equation}
  c_{\mathrm{r.m.s.}} = \sqrt{\langle c^2 \rangle}
\end{equation}

\subsubsection{Derivation Sketch}
\begin{enumerate}
  \item Consider one molecule of mass $m$, speed $c_x$ normal to a wall of area $A$
        in a box of length $L$.
  \item Momentum change per collision: $\Delta p = 2mc_x$.
  \item Time between successive collisions with the same wall: $\Delta t = 2L/c_x$.
  \item Force from one molecule: $F = \dfrac{2mc_x}{2L/c_x} = \dfrac{mc_x^2}{L}$.
  \item For $N$ molecules, extend to 3-D using isotropy:
        $\langle c_x^2 \rangle = \tfrac{1}{3}\langle c^2 \rangle$:
\end{enumerate}
\[
  p = \frac{Nm\langle c^2\rangle}{3V}
  \quad\Longrightarrow\quad
  pV = \tfrac{1}{3}Nm\langle c^2\rangle
\]

\subsection{Average Kinetic Energy of a Molecule}

Comparing $pV = NkT$ with $pV = \tfrac{1}{3}Nm\langle c^2\rangle$:

\begin{formulabox}{Average Translational Kinetic Energy}
\begin{equation}
  \langle E_k \rangle = \tfrac{1}{2}m\langle c^2\rangle = \tfrac{3}{2}kT
\end{equation}
The average kinetic energy of a molecule depends \textbf{only on temperature}.
\end{formulabox}

\begin{warnbox}
$\langle E_k \rangle = \tfrac{3}{2}kT$ gives the kinetic energy of \textbf{one
molecule}. For the total kinetic energy of $n$ moles use
$E_{\mathrm{total}} = \tfrac{3}{2}nRT$.
Do not confuse $k$ (Boltzmann constant, per molecule) with $R$ (molar gas constant,
per mole).
\end{warnbox}

% =============================================================================
\section{Topic 16 — Thermodynamics}
% =============================================================================

\subsection{Internal Energy}

\begin{defnbox}{Internal Energy $U$}
The \textbf{internal energy} of a system is the sum of the random kinetic and potential
energies of all its molecules. It is determined solely by the \textbf{state} of the
system (its temperature, pressure, volume and composition).
\end{defnbox}

For an \textbf{ideal gas} there are no intermolecular forces, so internal energy
consists entirely of kinetic energy:
\begin{equation}
  U = \tfrac{3}{2}NkT = \tfrac{3}{2}nRT \qquad \text{(ideal monatomic gas)}
\end{equation}
Hence for an ideal gas, $U$ depends \emph{only} on $T$.

For \textbf{real substances} the potential energy between molecules is also significant.
A rise in temperature increases kinetic energy; a change of state at constant
temperature changes potential energy only.

\subsection{Work Done by a Gas}

\begin{formulabox}{Work Done at Constant Pressure}
\begin{equation}
  W = p\,\Delta V
\end{equation}
\begin{tabular}{ll}
$W$        & = work done \textbf{by} the gas (J) \\
$p$        & = pressure (Pa) \\
$\Delta V$ & = change in volume (m$^3$) \\
\end{tabular}
\end{formulabox}

Sign convention: expansion ($\Delta V > 0$) means the gas does positive work on the
surroundings; compression ($\Delta V < 0$) means the surroundings do positive work on
the gas.

\subsection{The First Law of Thermodynamics}

\begin{defnbox}{First Law of Thermodynamics}
The increase in internal energy of a system equals the energy supplied to the system
by heating \textit{plus} the work done \textit{on} the system:
\begin{equation}
  \Delta U = q + W
\end{equation}
\end{defnbox}

Sign conventions (CIE syllabus form $\Delta U = q + W$):

\medskip
\begin{center}
\renewcommand{\arraystretch}{1.3}
\begin{tabular}{lll}
\toprule
\textbf{Symbol} & \textbf{Positive when\ldots} & \textbf{Negative when\ldots} \\
\midrule
$\Delta U$ & internal energy increases & internal energy decreases \\
$q$        & heat flows \textbf{into} system & heat flows \textbf{out} of system \\
$W$        & work done \textbf{on} system & work done \textbf{by} system \\
\bottomrule
\end{tabular}
\end{center}

\begin{tipbox}
CIE uses $\Delta U = q + W$ where $W$ is work done \textbf{on} the gas.
Some textbooks write $\Delta U = q - W$ where $W$ is work done \textbf{by} the gas.
Both are equivalent — always check your sign convention before substituting numbers.
\end{tipbox}

\subsubsection{Special Processes}

\begin{center}
\renewcommand{\arraystretch}{1.4}
\begin{tabular}{p{2.8cm}p{3.8cm}p{5.5cm}}
\toprule
\textbf{Process} & \textbf{Constraint} & \textbf{Consequence} \\
\midrule
Isothermal   & $\Delta T = 0 \Rightarrow \Delta U = 0$ & $q = -W$: heat in equals work out \\
Adiabatic    & $q = 0$                                 & $\Delta U = W$: compression raises $U$ \\
Isochoric    & $\Delta V = 0 \Rightarrow W = 0$        & $\Delta U = q$: all heating changes $U$ \\
Isobaric     & $p = \mathrm{const}$                    & $W = p\Delta V$; both $q$ and $W$ contribute \\
\bottomrule
\end{tabular}
\end{center}

% =============================================================================
\section{Formula Summary Sheet}
% =============================================================================

\begin{summarybox}
\renewcommand{\arraystretch}{1.6}
\begin{tabular}{p{6.5cm} p{5.5cm} p{1.8cm}}
\toprule
\textbf{Formula} & \textbf{Meaning} & \textbf{Topic} \\
\midrule
$T/\mathrm{K} = \theta/{}^{\circ}\mathrm{C} + 273.15$
  & Temperature conversion & 14 \\
$Q = mc\Delta T$
  & Heat and temperature change & 14 \\
$Q = mL$
  & Latent heat & 14 \\
$pV = nRT$
  & Ideal gas (molar form) & 15 \\
$pV = NkT$
  & Ideal gas (molecular form) & 15 \\
$k = R/N_A$
  & Boltzmann constant & 15 \\
$pV = \tfrac{1}{3}Nm\langle c^2\rangle$
  & Kinetic theory — pressure & 15 \\
$\langle E_k\rangle = \tfrac{1}{2}m\langle c^2\rangle = \tfrac{3}{2}kT$
  & Average KE per molecule & 15 \\
$c_{\mathrm{r.m.s.}} = \sqrt{\langle c^2\rangle}$
  & Root-mean-square speed & 15 \\
$U = \tfrac{3}{2}NkT = \tfrac{3}{2}nRT$
  & Internal energy (ideal gas) & 16 \\
$W = p\,\Delta V$
  & Work done by gas (const.\ $p$) & 16 \\
$\Delta U = q + W$
  & First law of thermodynamics & 16 \\
\bottomrule
\end{tabular}
\end{summarybox}

% =============================================================================
\section{Exam Technique and Problem-Solving Strategy}
% =============================================================================

\subsection{5-Step Strategy for Thermal Physics Problems}
\begin{enumerate}
  \item \textbf{Identify the topic}: temperature/heat, ideal gas, or thermodynamics?
  \item \textbf{List what is given} and what is being asked.
  \item \textbf{Convert units}: temperatures to K; volumes to m$^3$; pressures to Pa.
  \item \textbf{Select the correct formula} and rearrange \emph{before} substituting.
  \item \textbf{Check your answer}: correct units, sensible order of magnitude, correct sign.
\end{enumerate}

\subsection{Key Tips by Topic}

\subsubsection{Topic 14 — Temperature and Thermal Properties}
\begin{itemize}
  \item Split mixed heating problems into separate stages: warming + change of state
        + further warming.
  \item Remember $L_v \gg L_f$: for water,
        $L_v \approx 2.26\times10^6\,\mathrm{J\,kg}^{-1}$ vs
        $L_f \approx 3.34\times10^5\,\mathrm{J\,kg}^{-1}$.
  \item Calorimetry: set $IVt = mc\Delta T$; watch for heat losses in the experiment.
\end{itemize}

\subsubsection{Topic 15 — Ideal Gases and Kinetic Theory}
\begin{itemize}
  \item Use $pV = nRT$ for moles; use $pV = NkT$ for number of molecules.
  \item The kinetic theory derivation is examinable — practise the 3-step argument
        (momentum change $\to$ force $\to$ pressure, then isotropy to 3-D).
  \item R.m.s.\ speed depends on $\sqrt{T}$; doubling $T$ in kelvin multiplies
        $c_{\mathrm{r.m.s.}}$ by $\sqrt{2}$, \emph{not} by 2.
\end{itemize}

\subsubsection{Topic 16 — Thermodynamics}
\begin{itemize}
  \item Be rigorous about sign conventions: write them out explicitly before
        substituting.
  \item For an isothermal process with an ideal gas, $\Delta U = 0$ always.
  \item Work done \textbf{on} the gas appears as $+W$ in CIE's form of the first law.
\end{itemize}

% =============================================================================
\section{Worked Examples}
% =============================================================================

\begin{examplebox}{1 — Specific Heat Capacity}
\textbf{Question:} An electric heater of power 60\,W heats 0.50\,kg of water from
$20\,{}^{\circ}$C to $70\,{}^{\circ}$C in 175\,s. Calculate the specific heat capacity
of water.

\medskip\textbf{Solution:}

Energy supplied: $Q = Pt = 60 \times 175 = 1.05\times10^4\,\mathrm{J}$

Using $Q = mc\Delta T$:
\[
  c = \frac{Q}{m\,\Delta T} = \frac{1.05\times10^4}{0.50 \times 50}
    = 4200\,\mathrm{J\,kg^{-1}\,K^{-1}}
\]
\end{examplebox}

\begin{examplebox}{2 — Ideal Gas}
\textbf{Question:} A gas occupies $2.0\times10^{-3}\,\mathrm{m}^3$ at a pressure of
$1.5\times10^5\,\mathrm{Pa}$ and temperature $27\,{}^{\circ}$C. Calculate the number
of moles of gas.

\medskip\textbf{Solution:}

Convert: $T = 27 + 273.15 = 300.15 \approx 300\,\mathrm{K}$
\[
  n = \frac{pV}{RT}
    = \frac{1.5\times10^5 \times 2.0\times10^{-3}}{8.31 \times 300}
    = \frac{300}{2493}
    = 0.120\,\mathrm{mol}
\]
\end{examplebox}

\begin{examplebox}{3 — R.M.S.\ Speed}
\textbf{Question:} Calculate the r.m.s.\ speed of nitrogen molecules
($M_r = 28\,\mathrm{g\,mol}^{-1}$) at $20\,{}^{\circ}$C.

\medskip\textbf{Solution:}

$T = 293.15\,\mathrm{K}$;\quad
$m = \dfrac{28\times10^{-3}}{6.02\times10^{23}} = 4.65\times10^{-26}\,\mathrm{kg}$

\[
  \langle c^2\rangle = \frac{3kT}{m}
    = \frac{3 \times 1.38\times10^{-23} \times 293.15}{4.65\times10^{-26}}
    = 2.61\times10^5\,\mathrm{m^2\,s^{-2}}
\]
\[
  c_{\mathrm{r.m.s.}} = \sqrt{2.61\times10^5} = 511\,\mathrm{m\,s^{-1}}
\]
\end{examplebox}

\begin{examplebox}{4 — First Law of Thermodynamics}
\textbf{Question:} A gas is compressed. The work done \emph{on} the gas is 450\,J and
at the same time 200\,J of heat is removed from the gas. Find the change in internal
energy.

\medskip\textbf{Solution:}

Work done on gas: $W = +450\,\mathrm{J}$\quad (positive: on the gas)

Heat removed from gas: $q = -200\,\mathrm{J}$\quad (negative: leaves system)
\[
  \Delta U = q + W = -200 + 450 = +250\,\mathrm{J}
\]
Internal energy \emph{increases} by $250\,\mathrm{J}$.
\end{examplebox}

\newpage
% =============================================================================
\section{Practice Exam Questions}
% =============================================================================

\subsection*{Section A — Short Answer Questions}

\noindent\begin{minipage}{\textwidth}
\textbf{Q1.} Define specific heat capacity and state its SI unit.
\markalloc{2}
\end{minipage}

\bigskip
\noindent\begin{minipage}{\textwidth}
\textbf{Q2.} Distinguish between the specific latent heat of fusion and the specific
latent heat of vaporisation.
\markalloc{2}
\end{minipage}

\bigskip
\noindent\begin{minipage}{\textwidth}
\textbf{Q3.} State \emph{three} assumptions of the kinetic theory of gases.
\markalloc{3}
\end{minipage}

\bigskip
\noindent\begin{minipage}{\textwidth}
\textbf{Q4.} A 2.0\,kg block of ice at $0\,{}^{\circ}$C is melted and the resulting
water is heated to $100\,{}^{\circ}$C.
Given $L_f = 3.34\times10^5\,\mathrm{J\,kg}^{-1}$ and
$c_{\mathrm{water}} = 4200\,\mathrm{J\,kg}^{-1}\,\mathrm{K}^{-1}$, calculate the
total energy required.
\markalloc{3}
\end{minipage}

\bigskip
\noindent\begin{minipage}{\textwidth}
\textbf{Q5.} The temperature of an ideal gas is doubled on the kelvin scale (at
constant volume). State and explain what happens to:
\begin{itemize}[nosep]
  \item the r.m.s.\ speed of the molecules,
  \item the pressure of the gas.
\end{itemize}
\markalloc{4}
\end{minipage}

\subsection*{Section B — Longer Structured Questions}

\noindent\begin{minipage}{\textwidth}
\textbf{Q6.} An ideal gas is enclosed in a container of fixed volume
$5.0\times10^{-3}\,\mathrm{m}^3$. The gas contains $3.0\times10^{23}$ molecules at
a temperature of $300\,\mathrm{K}$.

\begin{enumerate}[label=(\alph*)]
  \item Calculate the pressure of the gas.
        \markalloc{2}
  \item The temperature is raised to $450\,\mathrm{K}$. Calculate the new pressure,
        assuming the gas remains ideal.
        \markalloc{2}
  \item Explain, using kinetic theory, why the pressure increases when the temperature
        is raised.
        \markalloc{3}
\end{enumerate}
\end{minipage}

\bigskip
\noindent\begin{minipage}{\textwidth}
\textbf{Q7.} A sample of gas absorbs $600\,\mathrm{J}$ of heat from its surroundings
and expands at a constant pressure of $1.2\times10^5\,\mathrm{Pa}$, increasing its
volume by $2.0\times10^{-3}\,\mathrm{m}^3$.

\begin{enumerate}[label=(\alph*)]
  \item Calculate the work done \emph{by} the gas.
        \markalloc{2}
  \item Using the first law of thermodynamics, calculate the change in internal energy
        of the gas.
        \markalloc{3}
  \item The gas is then compressed adiabatically back to its original volume. State
        what happens to the internal energy and explain why.
        \markalloc{2}
\end{enumerate}
\end{minipage}

\bigskip
\noindent\begin{minipage}{\textwidth}
\textbf{Q8.} Show that the average translational kinetic energy of a molecule of an
ideal gas is $\tfrac{3}{2}kT$.
\markalloc{3}
\end{minipage}

\newpage
% =============================================================================
\section{Mark Scheme and Answers}
% =============================================================================

\begin{markbox}
\textbf{Q1.} Specific heat capacity is the energy required to raise the temperature of
$1\,\mathrm{kg}$ of a substance by $1\,\mathrm{K}$ [1];\quad
unit: $\mathrm{J\,kg^{-1}\,K^{-1}}$ [1].

\medskip
\textbf{Q2.} Latent heat of \emph{fusion}: energy per unit mass to change solid to
liquid at constant temperature [1]. Latent heat of \emph{vaporisation}: energy per
unit mass to change liquid to gas at constant temperature [1].

\medskip
\textbf{Q3.} Any three from: large number of molecules [1]; random motion [1]; volume
of molecules negligible compared to container [1]; intermolecular forces negligible
between collisions [1]; perfectly elastic collisions [1]; collision duration negligible
[1]. \quad(3 marks total)

\medskip
\textbf{Q4.}
Melting: $Q_1 = 2.0 \times 3.34\times10^5 = 6.68\times10^5\,\mathrm{J}$ [1];
heating water: $Q_2 = 2.0 \times 4200 \times 100 = 8.40\times10^5\,\mathrm{J}$ [1];
total: $Q = 1.508\times10^6\,\mathrm{J} \approx 1.51\times10^6\,\mathrm{J}$ [1].

\medskip
\textbf{Q5.}
R.m.s.\ speed: $c_{\mathrm{r.m.s.}} \propto \sqrt{T}$, so doubling $T$ increases
$c_{\mathrm{r.m.s.}}$ by factor $\sqrt{2} \approx 1.41$ [1]; molecules have greater
kinetic energy [1].
Pressure: $p \propto T$ at constant $V$, so pressure doubles [1]; molecules collide
with walls more frequently and with greater momentum [1].

\medskip
\textbf{Q6(a).}
\[
  p = \frac{NkT}{V}
    = \frac{3.0\times10^{23} \times 1.38\times10^{-23} \times 300}{5.0\times10^{-3}}
    = 248\,\mathrm{Pa} \quad [2]
\]

\textbf{Q6(b).}
$p \propto T$ at constant $V$:\quad
$p_2 = 248 \times \dfrac{450}{300} = 372\,\mathrm{Pa}$ [2].

\textbf{Q6(c).}
At higher temperature molecules move faster [1]; they collide with the walls more
frequently [1]; each collision transfers greater momentum to the wall, so force and
pressure increase [1].

\medskip
\textbf{Q7(a).}
$W_{\text{by gas}} = p\Delta V = 1.2\times10^5 \times 2.0\times10^{-3}
= 240\,\mathrm{J}$ [2].

\textbf{Q7(b).}
Work done \emph{on} gas: $W = -240\,\mathrm{J}$ [1];\quad
$\Delta U = q + W = 600 + (-240) = +360\,\mathrm{J}$ [1];\quad
internal energy increases by $360\,\mathrm{J}$ [1].

\textbf{Q7(c).}
Internal energy increases [1]; no heat enters or leaves (adiabatic), so all work done
on the gas during compression goes entirely into increasing internal energy [1].

\medskip
\textbf{Q8.}
From kinetic theory: $pV = \tfrac{1}{3}Nm\langle c^2\rangle$ [1];
from ideal gas equation: $pV = NkT$;
equating: $\tfrac{1}{3}Nm\langle c^2\rangle = NkT$,
so $\tfrac{1}{2}m\langle c^2\rangle = \tfrac{3}{2}kT$ [1];
this is the average translational kinetic energy of one molecule [1].
\end{markbox}

\newpage
% =============================================================================
\section{Revision Checklist}
% =============================================================================

\begin{checklistbox}
\small
\renewcommand{\arraystretch}{1.5}
\begin{tabular}{p{11cm} >{\centering\arraybackslash}p{2.5cm}}
\toprule
\textbf{Learning Objective} & \textbf{Confidence (1--3)} \\
\midrule
\multicolumn{2}{l}{\textbf{Topic 14 — Temperature}} \\
$\square$\; Explain thermal equilibrium and state the zeroth law & \\
$\square$\; Convert temperatures between ${}^{\circ}$C and K using $T = \theta + 273.15$ & \\
$\square$\; State and explain examples of thermometric properties & \\
$\square$\; Define and use specific heat capacity: $Q = mc\Delta T$ & \\
$\square$\; Define specific latent heat; distinguish $L_f$ and $L_v$ & \\
$\square$\; Apply $Q = mL$; split calculations across stages correctly & \\
\midrule
\multicolumn{2}{l}{\textbf{Topic 15 — Ideal Gases}} \\
$\square$\; Use molar quantities and the Avogadro constant $N_A$ & \\
$\square$\; Recall and apply $pV = nRT$ and $pV = NkT$ & \\
$\square$\; State the relationship $k = R/N_A$ & \\
$\square$\; State all six assumptions of the kinetic theory & \\
$\square$\; Reproduce the derivation of $pV = \tfrac{1}{3}Nm\langle c^2\rangle$ & \\
$\square$\; Derive the result $\langle E_k\rangle = \tfrac{3}{2}kT$ & \\
$\square$\; Calculate r.m.s.\ speed from temperature and molar mass & \\
\midrule
\multicolumn{2}{l}{\textbf{Topic 16 — Thermodynamics}} \\
$\square$\; Explain internal energy in terms of molecular KE and PE & \\
$\square$\; State that internal energy of an ideal gas depends only on $T$ & \\
$\square$\; Calculate work done by/on a gas: $W = p\Delta V$ & \\
$\square$\; Apply the first law $\Delta U = q + W$ with correct sign conventions & \\
$\square$\; Analyse isothermal, adiabatic, isochoric and isobaric processes & \\
\bottomrule
\end{tabular}

\medskip
\begin{center}
\textit{Key:\quad 1 = Need more work\quad 2 = Getting there\quad 3 = Confident}
\end{center}
\end{checklistbox}

\vspace{1cm}
\begin{center}
{\large\bfseries\color{cieprimary} Good luck with your revision!}\\[0.4cm]
\textit{Remember: understanding \emph{why} formulae work is more powerful than
memorising them.}\\
\textit{Practise deriving $pV = \tfrac{1}{3}Nm\langle c^2\rangle$ and
$\langle E_k\rangle = \tfrac{3}{2}kT$ from first principles.}
\end{center}

\end{document}
